\pdfminorversion=7%
\documentclass[aspectratio=169,mathserif,notheorems]{beamer}%
%
\xdef\bookbaseDir{../../bookbase}%
\xdef\sharedDir{../../shared}%
\RequirePackage{\bookbaseDir/styles/slides}%
\RequirePackage{\sharedDir/styles/styles}%
\toggleToGerman%
%
\subtitle{10.~Fabrik-Datenbank:~Tabelle \emph{demand}}%
%
\begin{document}%
%
\startPresentation%
%
\section{Einleitung}%
%
\begin{frame}%
\frametitle{Einleitung}%
\begin{itemize}%
\item Jetzt können wir sowohl Kunden als auch Produkte in unserer Datenbank speichern.%
\item<2-> Als nächstes wollen wir auch Bestellungen speichern können.%
\item<3-> Wir wählen dafür ein extrem simplifiziertes Konzept\only<-3>{.}\uncover<4->{ %
Ein Kunde kann mit einer Bestellung eine Menge von genau einem Produkt ordern.%
}%
\item<5-> In echten Online-Shops können Bestellungen aus mehreren Produkten bestehen, Zahlungsoptionen haben, und vielleicht auch Lieferadressen\dots\uncover<6->{ {\dots}wir wollen unser Beispiel aber schön einfach halten.}%
\end{itemize}%
\end{frame}%
%
%
\section{Erstellen der Tabelle}%
%
\begin{frame}[t]%
\frametitle{Grobstruktur}%
\begin{itemize}%
\item Wir nennen unsere Tabelle für Bestellungen~\textil{demand}.\only<2-7>{ %
Ich würde sie lieber \sqlil{order} nennen, aber das ist ja ein SQL-Schlüsselwort\dots%
\item<3-> Die Tabelle muss für jede Bestellung folgende Daten speichern\uncover<4->{:%
\begin{itemize}%
\item den Kunde, der die Bestellung aufgibt\uncover<5->{,}%
\item<5-> das bestellte Produkt\uncover<6->{,}%
\item<6-> die Anzahl der bestellten Einheiten des Produktes\uncover<7->{, und}%
\item<7-> und das Datum, an dem die Bestellung aufgegeben wurde.%
\end{itemize}%
}}%
\only<-8>{\item<8-> Insgesamt könnten die Daten also wie folgt aussehen:}%
\item<9-> In den Spalten für Kunde und Produkt werden die Primärschlüssel aus den entsprechenden Tabellen stehen\dots%
\end{itemize}%
\only<2>{%
\bestPractice{neverUseKeywordsAsNames}{%
Never use \sql\ keywords or reserved words as names, e.g., for columns or tables\cite{B2025DS:SBPASG}.%
}}%%
\locateGraphic{9}{width=0.75\paperwidth}{graphics/dbStructureDemand}{0.125}{0.3}%
\end{frame}%
%
%
\begin{frame}[fragile,t]
\frametitle{Erstellen der Tabelle mit \texttt{CREATE TABLE}}%
\begin{itemize}%
\item Wir wissen ja schon, dass man Tabellen mit \sqlil{CREATE TABLE} erstellt.%
\item<4-> Wir haben dieses Kommando via \psql\ an den \postgresql-Server geschickt und das Erstellen der Tabelle hat funktioniert.%
\item<5-> Wieder gibt es einige interessante neue Komponenten\only<-5>{.}\uncover<6->{:~Was bedeutet \sqlil{REFERENCES}?\uncover<7->{ Was ist ein \sqlil{DATE}?\uncover<8->{ Schauen wir uns das Schritt für Schritt an.}}}%
\end{itemize}%
%
\gitLoadAndExecSQL{factory:create_table_demand}{}{factory}{create_table_demand.sql}{factory}{boss}{superboss123}%
\listingSQL{2}{factory:create_table_demand}{0.075}{0.18}{0.9}{0.84}
\listingOutput{3}{factory:create_table_demand}{,style=tool_style_sql}{0.05}{0.18}{0.9}{0.84}
\listingSQL{4-}{factory:create_table_demand}{0.17}{0.47}{0.66}{0.499}
\end{frame}
%
\begin{frame}%
\frametitle{Die Spalte \texttt{id}}%
\begin{itemize}%
\item Die erste Spalte nennen wir wieder \sqlil{id}.%
\item<2-> Sie hat den Typ \sqlil{INT GENERATED BY DEFAULT AS IDENTITY PRIMARY KEY}.%
\item<3-> In anderen Worten\uncover<4->{:%
\begin{itemize}%
\item ihre Werte sind ganzzahlig~(\sqlil{INT})\uncover<5->{,}%
\item<5-> werden automatisch von einer automatisch inkrementierten Sequenz generiert~(\sqlil{GENERATED BY DEFAULT AS IDENTITY})\uncover<6->{, und}%
\item<6-> stellen den Primärschlüssel für unsere Tabelle dar.%
\end{itemize}%
}%
\item<7-> Genaugenommen ist die Spalte wieder ein sogenannter \emph{Surrogate Key}, zu Deutsch ein Ersatzschlüssel.%
\end{itemize}%
\end{frame}%
%
\begin{frame}%
\frametitle{Die Spalte \texttt{customer}}%
\begin{itemize}%
\item Zu jeder Bestellung gehört ein Kunde.%
\item<2-> In der zweiten Spalte, die wir \sqlil{customer} nennen, wollen wir also die Kunden-Einträge referenzieren.%
\item<3-> Die Daten der Kunden sind in der Tabelle~\sqlil{customer} gespeichert.%
\item<4-> Sie werden über den Primärschlüssel identifiziert, den wir~\sqlil{id} genannt haben.%
\item<5-> Die Werte für \sqlil{id} sind einzigartig in der Tabelle~\sqlil{customer}:~Jeder Record $=$ jede Zeile der Tabelle hat einen anderen \sqlil{id}-Wert.%
\item<6-> Wir könnten also für jede Bestellung den \sqlil{id}-Wert des Kunden-Eintrags aus der \sqlil{customer}-Tabelle speichern.%
\item<7-> Mit diesem Wert können wir dann den Kunden-Record in der Tabelle \sqlil{customer} finden und z.B.~den Namen und die Adresse auslesen.%
\end{itemize}%
\end{frame}%
%
\begin{frame}[fragile,t]
\frametitle{Die Spalte \texttt{customer}}%
\begin{itemize}%
\item \alert<7>{Zu jeder Bestellung gehört ein Kunde}.%
\item Wir könnten für jede Bestellung den \sqlil{id}-Wert des Kunden-Eintrags aus der \sqlil{customer}-Tabelle\only<6->{ als Fremdschlüssel~(Foreign Key)} speichern\only<6->{\cite{PGDG:PD:FK}}.%
\only<-7>{%
\item Mit diesem Wert können wir dann den Kunden-Record in der Tabelle \sqlil{customer} finden und z.B.~den Namen und die Adresse auslesen.%
\item<2-> Wir würden also eine neue Spalte namens~\sqlil{customer} in der Tabelle~\sqlil{demand} erstellen, die vom Datentype~\sqlil{INT} ist~(denn die Spalte \sqlil{id} in er Tabelle~\sqlil{customer} ist ja auch ein~\sqlil{INT}).%
\item<3-> Wie stellen wir sicher, dass die Zahlen, die wir in der Spalte~\sqlil{customer} unserer neuen Tabelle~\sqlil{demand} speichern, wirklich passende Kunden-\sqlil{id}s sind?%
\item<4-> Irgendwie müssen wir dem DBMS sagen:~\emph{\inQuotes{Diese neue Spalte hier referenziert die Spalte~\sqlil{id} der Tabelle~\sqlil{customer}. Es dürfen nur Werte auftreten, die auch dort vorkommen!}}%
\item<5-> Das geht, in dem wir \sqlil{REFERENCES customer(id)} schreiben, wobei sich \sqlil{customer} auf die Tabelle~\sqlil{customer} und \sqlil{id}~auf die Spalte~\sqlil{id} der Tabelle~\sqlil{customer} bezieht.%
\item<6-> Durch \sqlil{REFERENCES customer(id)} wird die neue Spalte zum sogenannten Fremdschlüssel~(Foreign Key).%
\item<7-> Zusätzlich muss die Spalte als \sqlil{NOT NULL} markiert werden.%
}%
\end{itemize}%
\listingSQL{8}{factory:create_table_demand}{0.175}{0.34}{0.65}{0.63}
\end{frame}
%
\begin{frame}[fragile,t]
\frametitle{Die Spalte \texttt{product}}%
\begin{itemize}%
\item Zu jeder Bestellung gehört ein Produkt.%
\item<2-> Wir könnten für jede Bestellung den \sqlil{id}-Wert des Produkt-Eintrags aus der \sqlil{product}-Tabelle als Fremdschlüssel~(Foreign Key) speichern\cite{PGDG:PD:FK}.%
\only<-7>{%
\item<3-> Mit diesem Wert können wir dann den Produkt-Record in der Tabelle \sqlil{product} finden und z.B.~den Preis auslesen.%
\item<4-> Wir würden also eine neue Spalte namens~\sqlil{product} in der Tabelle~\sqlil{demand} erstellen, die vom Datentype~\sqlil{INT} ist~(denn die Spalte \sqlil{id} in er Tabelle~\sqlil{product} ist ja auch ein~\sqlil{INT}).%
\item<5-> Dazu müssen wir \sqlil{REFERENCES product(id)} schreiben, wobei sich \sqlil{product} auf die Tabelle~\sqlil{product} und \sqlil{id}~auf die Spalte~\sqlil{id} der Tabelle~\sqlil{product} bezieht.%
\item<6-> Dadurch wird sichergestellt, dass es nur Bestellungen geben kann, deren \sqlil{product}-Wert zu einer Zeile in der Tabelle~\sqlil{product} passt.%
\item<7-> Zusätzlich muss die Spalte als \sqlil{NOT NULL} markiert werden.%
}%
\end{itemize}%
\listingSQL{8}{factory:create_table_demand}{0.175}{0.34}{0.65}{0.63}
\end{frame}
%
\begin{frame}[t,fragile]
\frametitle{Die Spalte \texttt{amount}}%
\begin{itemize}%
\item Nun wollen wir noch speichern, wie viele Einheiten der Kunde von dem Produkt bestellt hat.%
\item<2-> Wir erstellen also eine Spalte namens~\sqlil{amount} vom Datentyp~\sqlil{INT}.%
\item<3-> Die Menge muss immer angegeben werden, also definieren wir sie als~\sqlil{NOT NULL}.%
\only<-7>{%
\item<4-> Stellt das aber sicher, dass die Menge \inQuotes{richtig} ist?%
\item<5-> Nein. Wir könnten 0 als Menge eingeben. Oder~-5.%
\item<6-> Wir fügen also ein \sqlil{CONSTRAINT} hinzu, das verlangt \sqlil{(amount > 0) AND (amount < 1000000)}.%
\item<7-> Randnotiz:~Solche Einschränkungen wären auch für die Tabelle~\sqlil{product} sinnvoll gewesen, allerdings war diese Lehreinheit schon kompliziert genug\dots%
}%
\end{itemize}%
\listingSQL{8}{factory:create_table_demand}{0.175}{0.365}{0.65}{0.63}
\end{frame}
%
\begin{frame}[t,fragile]
\frametitle{Die Spalte \texttt{ordered}}%
\begin{itemize}%
\item Wir müssen auch das Bestelldatum speichern.%
\item<2-> Wir erstellen also eine Spalte namens~\sqlil{ordered}\only<-8>{~(da \sqlil{WHEN} und \sqlil{DATE} reservierte SQL-Schlüsselworte sind\cite{PGDG:PD:SKW})}.%
\only<-8>{\item<3-> Aber welchen Datentyp sollten wir benutzen?%
\item<4-> Zum Glück gibt es in SQL Datentypen für Datums und Zeitangaben.}%
\item<5-> Hier verwenden wir den Datentyp~\sqlil{DATE}\only<-5,9->{.}\only<6-8>{, %
der ein Datum nach unserem (dem Gregorianischen) Kalender darstellt\cite{PGDG:PD:HU,G1582IG}.}%
\only<-8>{%
\item<7-> Also Schreibweise wird das ISO-Standardformat~\inQuotes{\textil{YYYY-MM-DD}}\cite{ISO860112019} verwendet, wobei \textil{YYYY}~das Jahr in vier Ziffern, \textil{MM}~den Monat in zwei Zifern, und \textil{DD}~den Tag in zwei Ziffern darstellen.%
\item<8-> Als Sicherheitscheck definieren wir ein \sqlil{CONSTRAINT}, dass sicherstellt das nur Bestellungen nach dem 01.10.2024 und vor dem 01.01.3000 eingegeben werden können.}%
\item<10-> Damit ist unsere Tabelle \sqlil{demand} fertig.%
\end{itemize}%
\listingSQL{9}{factory:create_table_demand}{0.175}{0.354}{0.65}{0.63}
\end{frame}
%
%
\section{Einfügen von Daten}%
%
\begin{frame}[t]%
\frametitle{Daten Einfügen}%
\begin{itemize}%
\item Nun haben wir die Tabelle \sqlil{demand} erstellt, aber sie ist leer.%
\item<2-> Wir wollen nun folgende acht Bestellungen speichern:%
\end{itemize}%
%
\locateGraphic{3}{width=0.7\paperwidth}{graphics/dbStructureDemand}{0.15}{0.33}%
\end{frame}%
%
%
\begin{frame}[fragile,t]
\frametitle{Füllen der Tabelle \texttt{demand}}%
\gitLoadAndExecSQL{factory:insert_into_table_demand}{}{factory}{insert_into_table_demand.sql}{factory}{boss}{superboss123}%
\listingSQL{1}{factory:insert_into_table_demand}{0.05}{0.13}{0.9}{0.91}
\listingOutput{2}{factory:insert_into_table_demand}{,style=tool_style_sql}{0.05}{0.13}{0.9}{0.91}
\end{frame}
%
\begin{frame}[t,fragile]
\frametitle{Schutz vor Fehlern}%
\begin{itemize}%
\item Schützt das DBMS uns auch wirklich vor Fehlern?%%
\item<2-> \only<-2>{%
Was passiert, wenn wir eine falsche Produkt-ID verwenden?}\only<3->{%
Wir können keine falsche Produkt-ID vwerden.}%
\item<4-> \only<-4>{%
Was passiert, wenn wir ein falsches Datum eingeben?}\only<5->{%
Wir können kein ungültiges Datum eingeben.}%
\end{itemize}%
%
\gitLoadAndExecSQL{factory:insert_into_table_demand_error_1}{}{factory}{insert_into_table_demand_error_1.sql}{factory}{boss}{superboss123}%
\listingSQL{2}{factory:insert_into_table_demand_error_1}{0.05}{0.33}{0.9}{0.665}
\listingSQLandOutput{3}{factory:insert_into_table_demand_error_1}{}{0.05}{0.33}{0.9}{0.665}
%
\gitLoadAndExecSQL{factory:insert_into_table_demand_error_2}{}{factory}{insert_into_table_demand_error_2.sql}{factory}{boss}{superboss123}%
\listingSQL{4}{factory:insert_into_table_demand_error_2}{0.05}{0.33}{0.9}{0.665}
\listingSQLandOutput{5-}{factory:insert_into_table_demand_error_2}{}{0.05}{0.33}{0.9}{0.665}
\end{frame}
%
\begin{frame}[t,fragile]
\frametitle{Understanding the Data}%
\begin{itemize}%
\item Wenn wir uns die Daten in unserer skizzierten Tabelle anschauen, dann können wir sie ganz gut verstehen, weil wir die Bedeutungen der IDs mit rangeschrieben haben.%
\item<2-> Ein \sqlil{SELECT} gibt uns aber nur die IDs für Produkte und Kunden {\dots} und ist daher schwer zu verstehen.%
\item<3-> In der nächsten Einheit lernen wir, wie wir die Daten aus den verschiedenen Tabellen zusammenführen und auswerten können.%
\end{itemize}%
\locateGraphic{1}{width=0.7\paperwidth}{graphics/dbStructureDemand}{0.15}{0.25}%
\listingOutput{2}{factory:insert_into_table_demand}{,style=tool_style_sql}{0.05}{0.34}{0.9}{0.65}
\end{frame}

\section{Zusammenfassung}%
%
\begin{frame}%
\frametitle{Zusammenfassung}%
\begin{itemize}%
\item Nun haben wir die dritte und letzte Tabelle unserer Fabrikdatenbank erstellt.%
\item<2-> Anders als die beiden Tabellen davor existiert sie nicht \inQuotes{alleine für sich.}%
\item<3-> Stattdessen referenziert sie die Tabellen \sqlil{customer} und \sqlil{product} über Fremdschlüssel, also in dem sie Spalten hat, die Werte von deren Primärschlüssel speichern.%
\item<4-> Somit können wir speichern, welcher Kunde welches Produkt bestellt hat.%
\item<5-> Als nächstes werden wir lernen, wie wir die Daten in unserer Datenbank auswerten können.%
\end{itemize}%
\end{frame}%
%
\endPresentation%
\end{document}%%
\endinput%
%
